% =============================================================================
%  INTRODUCCIÓN
% =============================================================================
\section{Introducción}

La Inteligencia Artificial (IA) avanza rápidamente, realizando cada vez más tareas que antes se otorgaban a una persona. Las empresas líderes, movidas por la competitividad y la ambición, buscan maximizar sus ganancias minimizando costos; reemplazar a un ser humano por un sistema que opere sin cansarse y sin parar les es simplemente mejor. Se suele decir que ``la inteligencia artificial es solo una herramienta que aumenta la productividad'', pero esta afirmación, cómoda y tranquilizadora, en un futuro podría convertirse en ``la herramienta que lo hace todo y mejor que todos'': la IA no es un fenómeno reciente, lleva décadas transformando silenciosamente distintos sectores, y su impacto actual no es comparable al de ninguna revolución tecnológica anterior en velocidad ni en alcance.

Como Ingeniero de Sistemas en entrenamiento, esta realidad no me resulta extraña. Los ingenieros somos aquellos encargados de construir estos sistemas, quienes escribimos los algoritmos, quienes diseñamos las arquitecturas que determinarán qué tan rápido y qué tan profundo llegará este cambio. Estamos construyendo algo que tiene el potencial de acabar con el sustento de muchísimas personas ---incluso con el nuestro propio--- y sin embargo, pocas veces en nuestra formación nos detenemos a preguntarnos si deberíamos. El Nuevo Humanismo, entendido en este curso como la búsqueda de un equilibrio entre el progreso tecnológico y el florecimiento genuino del ser humano para ayudar a la sociedad, nos exige exactamente esa pausa. En este ensayo sostengo que el avance de la Inteligencia Artificial, en el contexto de la globalización y el mercado utilitario, profundiza la desigualdad social y amenaza el propósito que muchas personas construyen a través del trabajo; y que la única respuesta coherente desde el Nuevo Humanismo es una ingeniería que se construya desde la consciencia ética, no desde la eficiencia ciega. Quiero recalcar, además, que el propósito y sentido de vida de un ser humano no yace en su utilidad.
